% appendix
\chapter{Filtering lists}
\label{sec:probe-filters}

\lstdefinelanguage{JavaScript}{
  keywords={break, case, catch, continue, debugger, default, delete, do, else, finally, for, function, if, in, const, instanceof, new, return, switch, this, throw, try, typeof, var, void, while, with},
  morecomment=[l]{//},
  morecomment=[s]{/*}{*/},
  morestring=[b]',
  morestring=[b]",
  sensitive=true
}



\section{Probe Filtering Resources}
\label{app:probe_filters}

This appendix documents the five external annotation resources adopted for technical probe filtering across all three datasets.
Each
resource targets distinct classes of measurement artefacts; together they define
the union-based exclusion criterion applied in
Section~\ref{subsec:step2}. Table~\ref{tab:filter_categories} summarises the
technical categories covered by each resource.

\subsection{Naeem \textit{et al.} (2014)}
\label{app:naeem2014}

Naeem \textit{et al.}~\cite{naeem2014} proposed a structured filtering framework
for the HumanMethylation450 array, designed to reduce false discoveries by
sequentially discarding probes on the basis of hybridisation specificity and
genomic integrity. The framework excludes probes
that multi-map or cross-hybridise to additional genomic loci, probes overlapping
repetitive elements (LINE, SINE, Alu), and probes overlapping regions harbouring insertion–deletion polymorphisms (INDELs). For SNP-affected probes, a
distinction is drawn between variants at the interrogated CpG or single-base
extension site---which directly compromise the methylation measurement and
are therefore excluded---and variants located nearby but not interfering in
bisulfite space, which are retained. The resulting exclusion set thus removes
probes affected by cross-reactivity, structural polymorphisms, and
measurement-critical SNPs.

\subsection{Chen \textit{et al.} (2013)}
\label{app:chen2013}

Chen \textit{et al.}~\cite{chen2013} empirically identified probes in the 450K
array exhibiting off-target hybridisation or overlap with common SNPs. In the
present work, only the cross-reactive probe list is used. This catalogue
enumerates approximately 29{,}000 multi-mapping probes whose methylation signal
cannot be unambiguously attributed to a single genomic locus, and whose inclusion
would therefore introduce spatially unresolved signal into the feature matrix.

\begin{table}[t]
\centering
\caption{Technical artefact categories covered by each filtering resource.}
\label{tab:filter_categories}
\scriptsize
\begin{tabular}{l c c c c c}
\toprule
\textbf{Category} &
\textbf{Naeem} &
\textbf{Chen} &
\textbf{Pidsley} &
\textbf{Zhou} &
\textbf{McCartney} \\
\midrule
Cross-hybridisation &
Yes &
Yes &
Yes &
\texttt{MASK\_mapping} &
Yes \\
SNP at CpG &
Yes &
--- &
Yes &
\texttt{MASK\_snp5}, \texttt{MASK\_extBase} &
--- \\
Nearby SNP &
Conditional &
--- &
--- &
--- &
--- \\
Structural variant &
Yes &
--- &
--- &
--- &
--- \\
Non-CpG probes &
--- &
--- &
Yes &
\texttt{MASK\_nonCG} &
Yes \\
\bottomrule
\end{tabular}
\end{table}


\subsection{Pidsley \textit{et al.} (2016)}
\label{app:pidsley2016}

Pidsley \textit{et al.}~\cite{pidsley2016} provide a platform-level assessment
of the EPIC array, releasing annotated probe lists that flag loci where technical
bias arises from design- or genome-related sources. Three categories are
relevant here. First, CpG-targeting probes with sequence homology to additional
genomic loci, which yield non-unique hybridisation and potentially ambiguous
$\beta$ estimates. Second, non-CpG-targeting probes affected by analogous
off-target issues, excluded due to limited interpretability in CpG-centric
analyses. Third, probes overlapping common genetic variation at three critical
positions: the interrogated CpG itself (polymorphic CpG, directly affecting the
presence of the CpG dinucleotide); the single-base extension site of Type~I
probes (perturbing extension chemistry and dye-intensity ratios); and within the
probe body (reducing binding affinity or altering hybridisation kinetics for
common variants).

\subsection{Zhou \textit{et al.} (2016)}
\label{app:zhou2016}

Zhou \textit{et al.}~\cite{zhou2017} released a unified probe annotation for
both 450K and EPIC arrays encoding technical reliability through a set of binary
mask columns (\texttt{MASK\_*}). The masks applied in this work flag probes with
low or inconsistent mapping quality (\texttt{MASK\_mapping}); Type~I probes
carrying a SNP at the extension base that induces a colour-channel switch
(\texttt{MASK\_typeINextBaseSwitch}); probes whose extension base is inconsistent
with the expected channel or CpG context (\texttt{MASK\_extBase}); probes with
non-unique 30-bp 3$'$ subsequences susceptible to cross-hybridisation
(\texttt{MASK\_sub30.copy}); probes overlapping an SNP within $\pm$5\,bp of the
interrogated CpG (\texttt{MASK\_snp5.common}); and probes overlapping SNPs with
global minor allele frequency above 1\% (\texttt{MASK\_snp5.GMAF1p}). The
composite \texttt{MASK\_general} field, which integrates mapping, SNP, and
cross-reactivity criteria, is used as the primary filter for general-purpose
exclusion. Probes overlapping RepeatMasker regions (\texttt{MASK\_rmsk15}) are
not excluded, consistent with standard practice.

\subsection{McCartney \textit{et al.} (2016)}
\label{app:mccartney2016}

McCartney \textit{et al.}~\cite{mccartney2016} assessed probe design artefacts
on the EPIC array and published supplementary tables enumerating
cross-hybridising CpG-targeting probes and cross-hybridising non-CpG-targeting
probes. These lists complement the Pidsley and Zhou resources by providing an
independent empirical characterisation of off-target hybridisation on the EPIC
platform.

The combined application of these resources ensures that downstream analyses operate on a probe set with maximised hybridisation specificity and minimal confounding by common genetic variation.